\section{Introduction}
\label{sec:intro}

Ask which floor of a tall building is the best place to be based --- the flat
you rent, the office you are assigned --- and the reflexive answer is ``the
middle''. The reasoning is the one-median on a line: if destinations are
spread uniformly over the floors, the median floor minimises the expected
distance to them. It is a clean argument, and it is wrong for most buildings,
because it treats the elevator as a distance. An elevator is a mechanism. When
you press the button the car is somewhere else --- at the lobby where the
controller parked it, or wherever it dropped its last passenger --- and it has
to come to you before it can take you anywhere. Most of your trips do not go
to a random floor at all: you leave the building and come back, so the lobby
is the destination or origin of the majority of them. Other people are using
the same car, and each of them adds a stop to your ride. And when you live on
the second floor you do not call the elevator to go down --- you take the
stairs, which is fast, whereas climbing back up is not.

This paper treats the question as a facility-location problem with a
mechanism inside it and answers it quantitatively. We write the expected trip
time of a resident on floor $h$ as a function of the lobby-trip fraction
$\alpha$, the elevator's kinematics (rated speed, acceleration, door and
transfer times), its idle parking policy (return to the lobby, stay put, park
in the middle, or --- the heavy-traffic limit --- be anywhere), the background
passenger rate, the number of cars, and the stairs. For one car and no traffic
the expectation has a closed form; for everything else we run a discrete-event
simulator with collective control and nearest-car dispatch, validated against
the closed form at zero traffic. The objective is the mean trip time, but we
also report the 90th percentile, the \emph{regret} of a floor (how much worse
than the best), and an \emph{indifference band} (the floors within 5\% of the
best), because the practical question is as much ``does it matter?'' as
``which floor?''.

Four findings organise the paper. (i) The middle-floor answer holds only in
the corner where no trip goes to the lobby and the car waits where you left
it; with lobby trips and lobby parking the optimum drops to the lowest
residential floor and the middle floor is far from indifferent. (ii)
Background traffic flattens the curve but does not move the optimum up until
the car is saturated, where an unexpected effect appears: on a loaded car the
low floors are the last to be served on the down sweep, so the optimum drifts
upward. (iii) Stairs shift the optimum upward by absorbing short trips, but
the asymmetry between walking down and climbing up, which we expected to be
the driver, contributes little. (iv) More cars make every floor faster and
make the choice of floor matter \emph{more}, not less, because what is left
of the trip is the ride, which is distance. A map over lobby fraction and load
turns these into a rule of thumb.
